\documentclass[11pt,a4paper]{article}
\usepackage[margin=24mm]{geometry}
\usepackage[T1]{fontenc}
\usepackage{lmodern,microtype,booktabs,graphicx,amsmath,amssymb,xcolor,fancyhdr,hyperref,caption,float}
\definecolor{navy}{HTML}{17324D}
\hypersetup{colorlinks=true,linkcolor=navy,urlcolor=navy}
\setlength{\parskip}{6pt}\setlength{\parindent}{0pt}
\setlength{\headheight}{15pt}\setlength{\emergencystretch}{2em}
\pagestyle{fancy}\fancyhf{}\fancyhead[L]{\small Multimodal Image--Text Fusion}\fancyhead[R]{\small Reproducible Experiment}\fancyfoot[C]{\thepage}
\newcommand{\fig}[3]{\begin{figure}[htbp]\centering\includegraphics[width=#2\linewidth]{figures/#1}\caption{#3}\end{figure}}
\begin{document}
\begin{titlepage}\centering\vspace*{30mm}{\Huge\bfseries\color{navy} Multimodal Fusion\par}\vspace{9mm}{\Large Reproducible Image--Text Classification\par}\vspace{8mm}{\large A controlled synthetic benchmark and technical evaluation\par}\vfill{\large Technical Project Report\par}\vspace{5mm}{\large October 2026\par}\vspace{20mm}{\small Dataset: 1,800 generated paired samples \quad | \quad 24$\times$24 images \quad | \quad 10 text features\par}\end{titlepage}
\tableofcontents\clearpage
\section{Executive summary}
This study tests whether joint learning from a noisy image and a noisy short text representation improves binary classification in a reproducible synthetic setting. Three independently trained networks share a fixed dataset partition. The observed held-out accuracies are \textbf{80.00\%} (image only), \textbf{77.11\%} (text only), and \textbf{100.00\%} (fusion). Zeroing the text representation at inference in the trained fusion model produces \textbf{47.78\%} accuracy. These numbers were computed by the accompanying executable and are not copied from the source proposal.

The synthetic task is deliberately constructed for cross-modal complementarity. Consequently, the results do not establish performance on natural-language, medical, or real-world images.
\section{Research objective and experimental design}
Each sample is $(X_i,T_i,y_i)$ with an image $X_i\in[0,1]^{1\times24\times24}$, a ten-dimensional binary bag-of-words vector $T_i\in\{0,1\}^{10}$, and a label $y_i\in\{0,1\}$. All 1,800 examples are produced using a local pseudorandom generator (seed 20261030). The partition is fixed at 1,260 train, 90 validation, and 450 test samples. Models are trained for exactly 12 epochs with Adam, learning rate $3\times10^{-3}$, mini-batch size 64, and multiclass cross-entropy. Validation is monitored but no test-dependent model selection is performed.
\subsection{Why the two modalities are complementary}
Two independent latent bits $a,b\in\{0,1\}$ define $y=a\oplus b$. The image encodes $a$ via the location of a bright square, while text encodes $b$ through a controlled word (\emph{opacity} versus \emph{involvement}). A separate oriented bar in the image is a noisy hint for $y$ (correct with probability 0.83); text orientation is independently a noisy hint (correct with probability 0.76). Hence neither latent bit alone determines the target, but both together do. These probabilities are \emph{generation parameters}, not reported empirical accuracies. This is a methodological demonstration, not a natural benchmark.
\fig{paired_examples.png}{0.94}{Actual examples sampled from the generated paired dataset. Text captions are generated from the same sample as each image.}
\section{Model architecture}
The image encoder uses convolution (1 to 8 channels, $3\times3$), ReLU, $2\times2$ max pooling, a second convolution (8 to 16 channels, $3\times3$), ReLU and pooling, then a 576-to-16 linear map with ReLU. The text encoder uses a 10-to-16 linear map and ReLU. Fusion concatenates the two 16-dimensional outputs, feeding a 32-to-24 ReLU layer and a 24-to-2 classifier. Unimodal baselines use the same 24-unit classifier head on their 16-dimensional embedding. Parameters are trained separately for each condition.
\begin{equation}h_x=f_\theta(X),\qquad h_t=g_\phi(T),\qquad z=W_2\operatorname{ReLU}(W_1[h_x;h_t]+b_1)+b_2.\end{equation}
\fig{architecture.png}{0.95}{Implemented two-branch neural network; dimensionalities match the source code.}
\section{Evaluation protocol and metrics}
Only the held-out 450-sample test split is used for final metrics. Accuracy is $(TP+TN)/N$; precision, recall and $F_1$ are macro-averaged over classes. The ablation is performed by replacing the learned 16-dimensional text embedding with zeros \emph{at inference}, without retraining. This tests reliance on the trained text pathway and is not identical to retraining an image-only model.
\begin{table}[htbp]\centering\caption{Measured held-out results from results/metrics.json.}\begin{tabular}{lrr}\toprule Condition & Accuracy (\%) & Macro $F_1$ \\ \midrule
Image-only & 80.00 & 0.799 \\
Text-only & 77.11 & 0.770 \\
Full fusion & 100.00 & 1.000 \\
Fusion, text embedding zeroed & 47.78 & 0.475 \\
\bottomrule\end{tabular}\end{table}
\fig{accuracy.png}{0.85}{Held-out accuracy computed from exact test-set predictions; all four conditions use the same test rows.}
\section{Learning dynamics}
The training trajectories below are logged epoch by epoch during the actual experiment. They are not illustrative or smoothed curves. A lower training loss does not on its own imply improved out-of-distribution robustness.
\fig{training_loss.png}{0.82}{Observed per-epoch mean training cross-entropy for each independently trained model.}
\section{Confusion analysis}
The trained fusion model correctly classifies 450 out of 450 held-out examples. Its confusion matrix (true labels are rows, predictions are columns) is
\begin{equation}\begin{bmatrix}207 & 0\\ 0 & 243\end{bmatrix}.\end{equation}
\fig{confusion_matrix.png}{0.68}{Fusion test predictions versus ground-truth class labels. Counts are directly calculated from saved predictions.}
\section{Latent representations}
The 32-dimensional learned fusion embeddings of the 450 held-out samples are projected using two-component PCA. The projection is a diagnostic display, not the classifier's decision boundary, and any apparent 2D separation must not be interpreted as proof of generalization.
\fig{latent_pca.png}{0.78}{PCA of the measured fusion embeddings for all held-out examples.}
\section{Interpretation and limitations}
Fusion exceeds the image-only baseline by 20.00 percentage points on this particular split. This is expected because the target contains an XOR relation across modalities: the image encodes one bit and text the other. A single held-out split and a single seed are insufficient for estimating uncertainty. Moreover, controlled keyword vectors are not general natural language; synthetic square-and-bar images are not photographs; and ablated inputs are out-of-distribution for the trained network. No external validity is claimed.
\section{Reproducibility and repository contents}
Run \texttt{python project.py --run} to regenerate all figures and JSON results. Run \texttt{python project.py --check} to verify stored metrics against the saved predictions and key split invariants. The report is compiled from \texttt{report.tex} using \texttt{pdflatex}. The experiment needs only Python and the pinned packages in \texttt{requirements.txt}; no network is needed once these packages are installed. The reference PDF is treated as a project specification; its example percentages and demonstrative curves are not presented as reproduced measurements.
\end{document}
